% ============================================================
% 第 21 章 直流电机的电枢绕组（教材第六篇，印刷页 235--243）
% 节名：第一节 直流电枢绕组的构成；第二节 单叠绕组；第三节 单波绕组
% ============================================================
\chapter{直流电机的电枢绕组}

直流电机的电枢绕组是产生电动势、通过电流产生转矩的核心部件。与交流绕组不同，它的线圈经换向片串联成一个\textbf{闭合回路}，再通过电刷分成若干并联支路。本章讨论两种基本型式：单叠绕组与单波绕组，重点掌握它们的\textbf{节距计算}与\textbf{并联支路数}。

\section{直流电枢绕组的构成}

\subsection{元件与换向片}
电枢绕组由许多\textbf{绕组元件}（线圈）组成，每个元件的两个端头分别接到两个相邻的换向片上。元件数 $S$、换向片数 $K$、虚槽数 $u$ 之间满足
\begin{equation}
S=K=uQ_a ,
\end{equation}
其中 $Q_a$ 为电枢槽数，$u$ 为每槽的虚槽数（每槽内并放的元件边对数）。为了分析方便，常把每个虚槽的上、下层元件边分别编号，用"虚槽"代替实际槽。

\subsection{节距}
直流电枢绕组有四个基本节距（均以\textbf{虚槽数}计）：
\begin{itemize}
  \item \textbf{第一节距} $y_1$：一个元件的两个元件边在电枢表面的距离，$y_1\approx\tau=\dfrac{Q_a}{2p}$（整距）；
  \item \textbf{第二节距} $y_2$：前一元件的下层边与后一元件的上层边之间的距离；
  \item \textbf{合成节距} $y$：相邻两元件对应边之间的距离，$y=y_1-y_2$（叠绕组）或 $y=y_1+y_2$（波绕组）；
  \item \textbf{换向器节距} $y_k$：相邻两元件所接换向片之间的距离，$y_k=y$。
\end{itemize}
为使 $y_1$ 取整数，$y_1$ 可整距、短距或长距，工程上多用整距或短距。

\subsection{并联支路数}
电枢绕组经电刷分成若干并联支路，设\textbf{支路对数}为 $a$，则电枢总电流 $I_a$ 与每条支路电流 $i_a$ 的关系为
\begin{equation}
I_a=2a\,i_a .
\end{equation}
$a$ 的大小由绕组型式决定，这是两种绕组最本质的区别。

\section{单叠绕组}

\subsection{连接规律}
单叠绕组的特点是\textbf{每个元件的末端接到相邻的下一个换向片上}，即
\begin{equation}
y=y_k=\pm1 .
\end{equation}
通常取 $y=+1$（右行绕组）。合成节距为
\begin{equation}
y=y_1-y_2=1\quad\Longrightarrow\quad y_2=y_1-1 .
\end{equation}
元件依次"叠"在一起前进，绕电枢一周后回到出发点附近，$p$ 对极电机需绕 $p$ 周才能把所有元件串完。

\subsection{并联支路数}
单叠绕组中，每个主磁极下的元件串联成一条支路，故
\begin{equation}
2a=2p\quad\Longrightarrow\quad a=p .
\end{equation}

\keyformula{\text{单叠绕组：}y=y_k=1,\qquad a=p,\qquad 2a=2p}

\subsection{特点}
\begin{itemize}
  \item 并联支路数多（$2a=2p$），每条支路电流小，适合\textbf{大电流、低电压}的电机；
  \item 电刷数目与极数相同，换向片间的电压较低；
  \item 常用于中小型直流电机及需要大电流的场合。
\end{itemize}

\section{单波绕组}

\subsection{连接规律}
单波绕组的特点是\textbf{每个元件的末端接到相隔较远的换向片上}，元件像波浪一样沿电枢前进，绕一周后回到出发点\textbf{相邻}的换向片，再继续第二周。其合成节距为
\begin{equation}
y=y_k=\frac{K\pm1}{p},
\end{equation}
取"$-$"号为左行绕组（常用），"$+$"号为右行绕组。合成节距与第一节距的关系为
\begin{equation}
y=y_1+y_2\quad\Longrightarrow\quad y_2=y-y_1 .
\end{equation}

\subsection{并联支路数}
单波绕组把所有同一极性下的元件都串联起来，只形成\textbf{两条}并联支路，与极数无关：
\begin{equation}
2a=2\quad\Longrightarrow\quad a=1 .
\end{equation}

\keyformula{\text{单波绕组：}y=y_k=\frac{K\pm1}{p},\qquad a=1,\qquad 2a=2}

\subsection{特点}
\begin{itemize}
  \item 并联支路数少（$2a=2$），每条支路串联元件多、电压高，适合\textbf{高电压、小电流}的电机；
  \item 同样容量下所需电刷数少（理论上只需两组电刷即可，实用上仍按极数布置以减小换向片间电压）；
  \item 常用于额定电压较高的直流电机及某些特殊场合（如电车牵引电机）。
\end{itemize}

\begin{center}
\begin{tabular}{@{}p{2.0cm}p{3.4cm}p{3.6cm}p{4.2cm}@{}}
\toprule
项目 & 单叠绕组 & 单波绕组 & 记忆要点 \\
\midrule
合成节距 & $y=y_k=\pm1$ & $y=y_k=\dfrac{K\pm1}{p}$ & 叠"1"、波"跨极" \\
$y_1$、$y_2$ 关系 & $y=y_1-y_2$ & $y=y_1+y_2$ & 叠减、波加 \\
支路对数 $a$ & $a=p$ & $a=1$ & 叠随极数、波恒为 1 \\
并联支路数 & $2a=2p$ & $2a=2$ & — \\
均压线 & 需要 & 一般不需要 & 由磁路不对称决定 \\
适用 & 大电流、低电压 & 高电压、小电流 & 由 $I_a=2ai_a$ 判断 \\
\bottomrule
\end{tabular}
\end{center}

\section{绕组连接步骤与均压线}

\subsection{连接的实际步骤}
以单叠绕组为例，画绕组展开图（或计算节距）的一般步骤如下：
\begin{enumerate}
  \item 由 $Q_a$、$u$、$p$ 求虚槽数 $Q=uQ_a$、元件数 $S=K=Q$、极距 $\tau=Q/(2p)$；
  \item 取第一节距 $y_1\le\tau$ 且为整数（整距或短距）；
  \item 由 $y=y_k=\pm1$ 定合成节距，再由 $y_2=y_1-y$ 求第二节距，检查 $y_2$ 是否为整数且合理；
  \item 从 1 号换向片（1 号元件上层边）出发，按 $+y_1$（上层边到下层边）、$-y_2$（下层边到下一元件上层边）逐元件编号，直到回到起点；
  \item 检查是否把所有元件都用上（单叠绕组需绕 $p$ 周），并确定电刷位置（放在几何中性线上，与主磁极轴线相差 $90^\circ$ 电角度处）。
\end{enumerate}
单波绕组的步骤相同，只是用 $y=(K\pm1)/p$、$y_2=y-y_1$，且绕行一周后回到起点\textbf{相邻}的换向片，继续第二周直到闭合。

\subsection{均压线}
单叠绕组中，若各极下磁路不对称（气隙不均、材料不均），同一极性各支路的感应电动势不相等，会在绕组内部产生\textbf{环流}，使损耗增大、发热不均。为此在电枢上装设\textbf{均压线}，把理论上等电位的换向片连接起来，让环流在均压线中流通而不经过电刷。

单波绕组的所有元件串联成两条支路，本身对磁路不对称不敏感，一般不需要均压线（这也是它在低速、多极电机中受欢迎的原因之一）。

\corepoint{两种绕组只需记两组式子：\textbf{单叠} $y=1$、$a=p$、$y_2=y_1-y$；\textbf{单波} $y=(K\pm1)/p$、$a=1$、$y_2=y-y_1$。判断"该选哪种"的标准是电流与电压：绕组并联支路越多，每条支路分担的电流越小、总电压越低——所以大电流选单叠、高电压选单波。另外注意 $y$、$y_1$、$y_2$、$y_k$ 都要取整数，$y_1$ 须满足 $y_1\le\tau$ 且为整数。}
